\section{Limitations and Discussion}

The control guide is neither a semantic oracle nor a complete characterization
of planning relevance.  It can exploit a shortcut, omit a static goal, or
produce a small gradient after its decoder has learned an important feature.
The hard rule consequently provides optimization priority, not a guarantee of
physical sufficiency.  Soft residual admission makes this uncertainty explicit
but may allow a nuisance to accumulate again over long training.

Our geometric and frozen-probe analyses establish representation structure and
decodability, not causal use.  A future-latent intervention that permutes
candidate action--outcome pairings inside CEM is needed to show that improved
predictions actually determine selected actions.  Likewise, current principal
results begin with seed zero and 50-episode checkpoint evaluations.  Final
claims require multiple training seeds, more evaluation episodes, and
predeclared checkpoint aggregation.

Finally, the observed watermark collapse is task dependent: tagged TwoRooms
and Cube baselines degrade less severely than tagged PushT.  This variation is
scientifically useful because it prevents the method from being presented as a
universal filter.  The broader claim is that predictive objectives can
misallocate a planner-facing representation, and that their encoder-side
optimization authority should be measured and controlled.

\section{Conclusion}

Predictive accuracy and representation utility are not interchangeable.  A
JEPA predictor may learn a visually predictable factor that is irrelevant to
future decisions, while a short-horizon control loss may fail to cover
information essential for long-horizon planning.  By separating predictor
learning from encoder allocation, \shortmethod{} provides a practical probe
and intervention at this boundary.  The resulting evidence supports a more
careful principle: use control signals to prioritize predictive updates, but
do not interpret instantaneous orthogonality as proof that an update is
useless.
